\section{Palindrome Partitioning II}
\label{sec:dp-palindrome-partitioning-2}

\problemstatement{Given a string $s$, partition it into the minimum
possible number of substrings such that every substring is a
palindrome. Return the \textbf{minimum number of cuts} needed
(one fewer than the number of pieces).}

\begin{patternbox}
This problem layers \textbf{two} DP tables, each from a different
chapter's family: an interval DP table (Section~\ref{sec:dp-partition-intro}'s
style) precomputing \emph{which substrings are palindromes}, feeding
into a 1D ending-here table (the DP on LIS chapter's style) computing
the minimum cuts.
\end{patternbox}

\subsection*{Step 1: Precomputing palindromic substrings (interval DP)}

Let $\textit{isPal}(i,j)$ be true exactly when $s[i..j]$ is a
palindrome. The classic interval recurrence:
\[
  \textit{isPal}(i,j) = \big(s[i]=s[j]\big) \land
  \big(j-i<2 \lor \textit{isPal}(i{+}1,j{-}1)\big).
\]
A single character ($i=j$) or two equal characters ($j=i+1$,
$s[i]=s[j]$) is trivially a palindrome; otherwise, $s[i..j]$ is a
palindrome exactly when its outer characters match \emph{and} the
inner substring $s[i{+}1..j{-}1]$ is itself a palindrome. This table is
filled by increasing range length, exactly like every other interval DP
in this chapter.

\subsection*{Step 2: Minimum cuts (1D ending-here DP)}

Let $\textit{cuts}[i]$ be the minimum cuts needed for the prefix
$s[0..i]$. For every $j$ from $0$ to $i$ where $s[j..i]$ is a
palindrome (looked up from Step 1's table in $O(1)$), this palindrome
can be the \textbf{last piece} of the partition, requiring whatever cuts
were needed for the prefix before it, plus one more cut to separate
this final piece (or zero cuts if $j=0$, since the whole prefix is one
single piece with no preceding pieces):
\[
  \textit{cuts}[i] = \min_{j\le i,\ \textit{isPal}(j,i)}
  \big(j=0\ ?\ 0 : \textit{cuts}[j{-}1] + 1\big).
\]

\subsection*{C++ implementation}

\begin{lstlisting}
#include <bits/stdc++.h>
using namespace std;

int minCutsPalindromePartition(string s) {
    int n = s.size();
    vector<vector<bool>> isPal(n, vector<bool>(n, false));

    for (int len = 1; len <= n; len++) {
        for (int i = 0; i + len - 1 < n; i++) {
            int j = i + len - 1;
            if (s[i] == s[j] && (len <= 2 || isPal[i + 1][j - 1])) {
                isPal[i][j] = true;
            }
        }
    }

    vector<int> cuts(n, INT_MAX);
    for (int i = 0; i < n; i++) {
        for (int j = 0; j <= i; j++) {
            if (isPal[j][i]) {
                int prevCuts = (j == 0) ? 0 : cuts[j - 1] + 1;
                cuts[i] = min(cuts[i], prevCuts);
            }
        }
    }
    return cuts[n - 1];
}

int main() {
    cout << minCutsPalindromePartition("aab") << endl; // 1
    return 0;
}
\end{lstlisting}

\subsection*{Dry run}

\dryrun{Take $s = \texttt{"aab"}$.}

Step 1: $\textit{isPal}$ is true for every single character
($(0,0),(1,1),(2,2)$), for $(0,1)$ (\texttt{"aa"}), and false for
$(1,2)$ (\texttt{"ab"}) and $(0,2)$ (\texttt{"aab"}, outer characters
$a,b$ mismatch).

Step 2: $\textit{cuts}[0]=0$ (single character \texttt{"a"}, no cuts
needed). $\textit{cuts}[1]$: $j=0$ gives \texttt{"aa"}, a palindrome,
contributing $0$; $j=1$ gives \texttt{"a"} alone, contributing
$\textit{cuts}[0]+1=1$. Minimum: $\textit{cuts}[1]=0$.
$\textit{cuts}[2]$: $j=2$ gives \texttt{"b"} alone, contributing
$\textit{cuts}[1]+1=1$; no other $j$ gives a palindrome ending at
index $2$. $\textit{cuts}[2]=1$. Answer: $1$ cut, partitioning
\texttt{"aa"} $\mid$ \texttt{"b"}.

\begin{complexitybox}
\textbf{Time Complexity.} $O(n^2)$ for Step 1 (interval DP, filled in
$O(1)$ per cell here since no further split-point search is needed --
only the immediately smaller inner range is checked), plus $O(n^2)$ for
Step 2 (an $O(n)$ scan per index): $O(n^2)$ overall.
\textbf{Space Complexity.} $O(n^2)$ for \textit{isPal}, plus $O(n)$ for
\textit{cuts}.
\end{complexitybox}

\begin{takeawaybox}
Not every interval DP needs the full $O(n)$ split-point search this
chapter opened with -- \textit{isPal} is an interval table whose
recurrence only ever looks at \emph{one} smaller inner range, not a
search over all of them, because a palindrome's defining property
(matching outer characters) only has one way to shrink. Recognizing
when an interval recurrence collapses to a single inner lookup (rather
than a full split-point scan) is what keeps this problem at $O(n^2)$
instead of the $O(n^3)$ every other section in this chapter needs.
\end{takeawaybox}
